% Generated from pinned audit evidence.
\newcommand{\TrainSeconds}{361.6}
\newcommand{\TrainMemory}{0.543}
\newcommand{\FirstIdentity}{1.17}
\newcommand{\FirstShuffled}{1.17}
\newcommand{\FirstZero}{0.00}
\newcommand{\BlockIdentity}{0.586}
\newcommand{\BlockGain}{0.098}
\newcommand{\BlockLow}{0.000}
\newcommand{\BlockHigh}{0.307}
\newcommand{\BlockFone}{0.320}
\newcommand{\RunDecision}{reject}

% Generated from the pinned V54 audit; no fitted presentation values.
\newcommand{\DensitySeconds}{207.8}
\newcommand{\DensityMemory}{0.400}
\newcommand{\DensityNll}{15.219}
\newcommand{\DensityShuffledNll}{15.254}
\newcommand{\DensityAdvantage}{0.0355}
\newcommand{\DensityLow}{-0.1014}
\newcommand{\DensityHigh}{0.1693}
\newcommand{\DensityIdentity}{1.562}
\newcommand{\DensityFone}{0.350}
\newcommand{\DensityCopyFone}{0.99949}
\newcommand{\DensityDecision}{reject}

% Generated from the locked V55 matched audit.
\newcommand{\JointNll}{7.026}
\newcommand{\ControlNll}{7.579}
\newcommand{\JointNllGain}{0.553}
\newcommand{\JointNllLow}{0.450}
\newcommand{\JointNllHigh}{0.665}
\newcommand{\JointTime}{342.4}
\newcommand{\ControlTime}{218.0}
\newcommand{\JointMemory}{0.648}
\newcommand{\ControlMemory}{0.400}
\newcommand{\JointSampleFone}{0.362}
\newcommand{\JointPrefixGain}{0.185}
\newcommand{\JointDecision}{reject}

% Generated from unchanged V56 report.
\newcommand{\BeamAccuracy}{11.621}
\newcommand{\BeamShuffled}{9.863}
\newcommand{\GreedyAccuracy}{8.252}
\newcommand{\BeamGain}{3.369}
\newcommand{\BeamGainLow}{2.351}
\newcommand{\BeamGainHigh}{4.412}
\newcommand{\BeamHistoryGain}{1.758}
\newcommand{\BeamHistoryLow}{0.966}
\newcommand{\BeamHistoryHigh}{2.552}
\newcommand{\BeamRepeat}{25.446}
\newcommand{\TargetRepeat}{0.112}
\newcommand{\BeamSeconds}{269.3}
\newcommand{\BeamMemory}{0.154}
\newcommand{\BeamDecision}{reject\_decoder}

% Generated from the completed V57 audit.
\newcommand{\QaExact}{3.906}
\newcommand{\QaFone}{15.718}
\newcommand{\QaBlankExact}{0.391}
\newcommand{\QaBlankFone}{5.553}
\newcommand{\QaControlExact}{0.000}
\newcommand{\QaControlFone}{6.582}
\newcommand{\QaQuestionGain}{10.165}
\newcommand{\QaQuestionLow}{7.043}
\newcommand{\QaQuestionHigh}{13.371}
\newcommand{\QaPairExact}{0.254}
\newcommand{\QaMotor}{100.000}
\newcommand{\QaSeconds}{278.572}
\newcommand{\QaMemory}{1.053}
\newcommand{\QaControlSeconds}{269.126}
\newcommand{\QaControlMemory}{1.049}
\newcommand{\QaDecision}{reject\_reading\_model}

% Generated from the fixed autonomous verification receipt.
\newcommand{\CodecLanguageFOne}{32.671}
\newcommand{\CodecLanguageBlankFOne}{19.508}
\newcommand{\CodecLanguageQuestionGain}{13.163}
\newcommand{\CodecLanguageCILow}{10.933}
\newcommand{\CodecLanguageCIHigh}{15.730}

\newcommand{\PaperTitle}{Learning conditional language\\in continuous image space}
\newcommand{\PaperSummary}{Written language includes historical, handwritten and transitional forms that a fixed character inventory cannot faithfully represent. We investigate independent language modelling with writing images as input and output on one consumer GPU. The actors consume and generate physical rasters without a character vocabulary. An earlier compact reader answers 25 of 256 development questions exactly, but only four of 393 shared-paragraph pairs have both answers correct. We then separate physical rendering from conditional answer selection. A frozen continuous writing codec reconstructs 128 development blocks at 99.998\% foreground F1, with 126 exact blocks. An inherited 7.539B-parameter actor learns to predict this output space by adapting 9.352M parameters for 10,000 updates on one RTX 4090 D, with 16.965 GiB peak allocated memory. Its autonomous Chinese answers reach only \CodecLanguageFOne\% F1 and zero exact matches out of 128, despite a \CodecLanguageQuestionGain-point effect of question pixels and correct stopping. A fixed 32-case diagnostic finds poor initial and exact-history predictions, including on training examples, while supplied-code reconstruction remains faithful. These results identify a conditional-prediction failure preceding generated-history drift. Useful general image-language generation and historical-form reasoning remain open objectives.}
